% !TEX program = xelatex
\documentclass[11pt,a4paper]{article}

\usepackage[UTF8]{ctex}
\usepackage[margin=2.2cm]{geometry}
\usepackage{amsmath, amssymb, amsthm}
\usepackage{listings}
\usepackage{xcolor}
\usepackage{tikz}
\usepackage{booktabs}
\usepackage{multirow}
\usepackage{hyperref}
\usetikzlibrary{positioning, arrows.meta, shapes.geometric, fit, calc, decorations.pathreplacing}

\hypersetup{
  colorlinks=true,
  linkcolor=blue!60!black,
  urlcolor=teal!70!black,
}

\lstdefinelanguage{Rust}{
  morekeywords={fn,let,mut,pub,struct,impl,use,mod,self,Self,Option,Some,None,return,if,else,for,while,loop,match,break,continue,as,where,trait,type,enum,move,ref,const,static,unsafe,extern,crate,super,dyn,async,await,box,in},
  sensitive=true,
  morecomment=[l]{//},
  morecomment=[s]{/*}{*/},
  morestring=[b]",
}

\lstset{
  language=Rust,
  basicstyle=\ttfamily\footnotesize,
  keywordstyle=\color{blue!60!black}\bfseries,
  commentstyle=\color{green!40!black}\itshape,
  stringstyle=\color{red!60!black},
  numbers=left,
  numberstyle=\tiny\color{gray},
  numbersep=8pt,
  frame=single,
  framerule=0.4pt,
  rulecolor=\color{gray!40},
  breaklines=true,
  breakatwhitespace=true,
  showstringspaces=false,
  tabsize=4,
  columns=fullflexible,
  keepspaces=true,
}

\title{\textbf{halfedge\,: 稀疏线性代数模块设计文档} \\ \large \texttt{src/linalg.rs}}
\author{halfedge 库}
\date{2026-07-05}

\begin{document}
\maketitle
\tableofcontents

\section{模块定位}

\texttt{linalg.rs} 是 \texttt{halfedge} 库的\textbf{稀疏线性代数工具}模块。
在 \texttt{sprs} 的 \texttt{CsMat} 稀疏矩阵之上提供三层能力：

\begin{itemize}
  \item \texttt{SparseSystem}：从网格拉普拉斯矩阵的三元组构建稀疏对称系统；
  \item \texttt{conjugate\_gradient}：共轭梯度法求解 $Ax = b$（$A$ 对称正定）；
  \item 辅助工具：正则化对角偏移 \texttt{regularize\_diagonal}、残差范数计算。
\end{itemize}

\textbf{使用场景}：
\begin{itemize}
  \item 参数化：求解调和 / LSCM / Tutte 线性系统；
  \item 测地线：Heat Method 中的两次 Poisson 求解；
  \item 共形映射：离散调和映射。
\end{itemize}

\begin{center}
\renewcommand{\arraystretch}{1.18}
\begin{tabular}{@{}lll@{}}
  \toprule
  \textbf{类别} & \textbf{API} & \textbf{语义} \\
  \midrule
  \multirow{2}{*}{构建器}
    & \texttt{SparseSystem::new(dim)}         & 新建 $dim \times dim$ 空稀疏系统 \\
    & \texttt{SparseSystem::add(i, j, val)}   & 对称累加 $A_{ij} \mathrel{+}= val$（同时写 $A_{ji}$） \\
  \midrule
  \multirow{3}{*}{构建器}
    & \texttt{SparseSystem::add\_diag(i, val)}& 仅累加对角元 $A_{ii} \mathrel{+}= val$ \\
    & \texttt{SparseSystem::finish(self)}     & 转为 CSR 格式 \texttt{CsMat} \\
    & \texttt{SparseSystem::dim(\&self)}      & 返回矩阵维度 \\
  \midrule
  求解器
    & \texttt{conjugate\_gradient(a, b, max\_iter, tol)} & 求解 $Ax=b$；未收敛返回 \texttt{None} \\
  \midrule
  辅助
    & \texttt{regularize\_diagonal(a, lambda)} & 对角正则化 $A \leftarrow A + \lambda I$ \\
  \bottomrule
\end{tabular}
\end{center}

\section{数据结构：SparseSystem}

\texttt{SparseSystem} 是一个\textbf{正在构建中的稀疏对称线性系统}。内部用
\texttt{sprs::TriMat<f64>}（三元组格式）收集矩阵元素，调用 \texttt{finish}
后转为 CSR 格式的 \texttt{CsMat}，可直接用于共轭梯度求解。

\begin{lstlisting}[language=Rust]
pub struct SparseSystem {
    tri: TriMat<f64>,
    dim: usize,
}
\end{lstlisting}

\subsection{对称累加语义}

\texttt{add(i, j, val)} 实现\textbf{对称累加}：当 $i \ne j$ 时同时写入
$A_{ij}$ 与 $A_{ji}$，保证矩阵对称；当 $i = j$ 时只写一次对角元：

\[
  A_{ij} \mathrel{+}= val,\quad
  A_{ji} \mathrel{+}= val \quad (i \ne j).
\]

\texttt{add\_diag(i, val)} 仅累加对角元 $A_{ii} \mathrel{+}= val$，适用于
拉普拉斯矩阵中顶点 \textbf{degree} 项的填充。

\subsection{索引越界守卫}

\texttt{add} 与 \texttt{add\_diag} 在写入前检查 $i, j < \texttt{dim}$，
越界索引静默丢弃（不 panic）。这是面向\textbf{网格拓扑}场景的容错策略：
遍历半边 / 顶点时偶有 \texttt{VertexId} 被删除或越界，应跳过而非中断构建。

\subsection{从 TriMat 到 CsMat}

\texttt{finish(self)} 调用 \texttt{TriMat::to\_csr()} 完成\textbf{一次性压缩}：
合并重复三元组、按行排序、丢弃显式零元。此后矩阵即不可变，仅供
\texttt{conjugate\_gradient} 读取。

\section{共轭梯度法}

\texttt{conjugate\_gradient(a, b, max\_iter, tol)} 求解对称正定（SPD）系统
$Ax = b$。算法复杂度 $O(\text{iter} \cdot \text{nnz})$，远优于直接求逆。

\subsection{算法公式}

\textbf{初始化}（取 $x_0 = 0$）：
\[
  r_0 = b - A x_0 = b,\qquad p_0 = r_0,\qquad \rho_0 = r_0 \cdot r_0.
\]

\textbf{迭代}（$k = 0, 1, \dots, \text{max\_iter}-1$）：
\begin{align}
  \alpha_k &= \frac{r_k \cdot r_k}{p_k \cdot A p_k}, \\
  x_{k+1}  &= x_k + \alpha_k\, p_k, \\
  r_{k+1}  &= r_k - \alpha_k\, A p_k, \\
  \beta_k  &= \frac{r_{k+1} \cdot r_{k+1}}{r_k \cdot r_k}, \\
  p_{k+1}  &= r_{k+1} + \beta_k\, p_k.
\end{align}

\textbf{收敛判据}（相对残差）：
\[
  \frac{\|r_k\|_2}{\|b\|_2} < \text{tol} \quad\Rightarrow\quad \text{返回 } \texttt{Some}(x_k).
\]

\textbf{未收敛}：若达到 \texttt{max\_iter} 仍未满足判据，返回 \texttt{None}，
由调用方决定是否增大迭代数或正则化。

\subsection{流程图}

\begin{center}
\resizebox{\textwidth}{!}{%
\begin{tikzpicture}[
  node distance=0.55cm,
  proc/.style={draw, rounded corners=2pt, minimum width=10cm, minimum height=0.7cm, align=center, font=\small},
  dec/.style={diamond, draw, aspect=2.6, fill=orange!12, inner sep=1pt, font=\small, align=center, minimum width=4cm},
  op/.style={-Stealth, thick},
  startend/.style={draw, rounded corners=10pt, minimum width=2.8cm, minimum height=0.7cm, font=\small, fill=gray!12},
  term/.style={draw, rounded corners=2pt, fill=green!12, minimum width=4cm, minimum height=0.7cm, font=\small, align=center}
]
  \node[startend] (s) {开始};
  \node[proc, below=of s, fill=blue!8] (p1) {初始化：$x_0 = 0$,\ $r_0 = b$,\ $p_0 = r_0$,\ $\rho_0 = r_0 \cdot r_0$};
  \node[proc, below=of p1, fill=blue!4] (p2) {计算 $A p_k$（稀疏矩阵-向量乘）};
  \node[proc, below=of p2, fill=blue!4] (p3) {$\alpha_k = \rho_k \,/\, (p_k \cdot A p_k)$};
  \node[proc, below=of p3, fill=blue!4] (p4) {$x_{k+1} = x_k + \alpha_k\, p_k$;\quad $r_{k+1} = r_k - \alpha_k\, A p_k$};
  \node[dec, below=of p4] (d1) {$\|r_{k+1}\| / \|b\| < \text{tol}$？};
  \node[term, right=1.6cm of d1] (t1) {返回 \texttt{Some}($x_{k+1}$)};
  \node[proc, below=of d1, fill=yellow!15] (p5) {$\beta_k = (r_{k+1}\!\cdot r_{k+1}) \,/\, \rho_k$；\ $p_{k+1} = r_{k+1} + \beta_k\, p_k$；\ $\rho_{k+1} \leftarrow \rho_k$};
  \node[dec, below=of p5] (d2) {已达 \texttt{max\_iter}？};
  \node[term, right=1.6cm of d2] (t2) {返回 \texttt{None}};
  \node[proc, below=of d2, fill=blue!4] (p6) {$k \leftarrow k+1$，回到计算 $A p_k$};
  \draw[op] (s) -- (p1);
  \draw[op] (p1) -- (p2);
  \draw[op] (p2) -- (p3);
  \draw[op] (p3) -- (p4);
  \draw[op] (p4) -- (d1);
  \draw[op] (d1) -- node[above, font=\scriptsize] {是} (t1);
  \draw[op] (d1) -- node[right, font=\scriptsize] {否} (p5);
  \draw[op] (p5) -- (d2);
  \draw[op] (d2) -- node[above, font=\scriptsize] {是} (t2);
  \draw[op] (d2) -- node[right, font=\scriptsize] {否} (p6);
  \draw[op] (p6.east) .. controls +(2.2,0) and +(2.2,0) .. (p2.east);
\end{tikzpicture}
}
\end{center}

\subsection{退化守卫}

\begin{itemize}
  \item $n = 0$（空系统）直接返回 \texttt{Some(Vec::new())}；
  \item $\|b\| < 10^{-30}$（右端近似零）返回零解 \texttt{vec![0.0; n]}，避免除零；
  \item 维度不匹配触发 \texttt{assert\_eq!}（编程错误，应 panic）。
\end{itemize}

\section{对角正则化}

\texttt{regularize\_diagonal(a, lambda)} 对稀疏 SPD 矩阵做对角偏移：
\[
  A_{\text{reg}} = A + \lambda\, I.
\]

\textbf{用途}：处理拉普拉斯矩阵的\textbf{零空间}。图拉普拉斯 $L$ 半正定，
常零空间维数 $\ge 1$（对应常数向量），CG 不收敛。加上 $\lambda I$ 后
$L + \lambda I$ 严格正定，CG 可收敛；典型 $\lambda \in [10^{-6}, 10^{-3}]$。

\textbf{实现细节}：遍历对角索引 $i \in [0, n)$，通过 \texttt{CsMat::get\_mut(i, i)}
就地累加。对拉普拉斯矩阵对角线一定非零（顶点 degree），故无需处理
\texttt{(i, i)} 显式零元的情况。

\section{退化守卫汇总}

\begin{center}
\renewcommand{\arraystretch}{1.18}
\begin{tabular}{@{}lll@{}}
  \toprule
  \textbf{API} & \textbf{退化场景} & \textbf{处理} \\
  \midrule
  \texttt{SparseSystem::add}        & $i, j \ge \text{dim}$（越界） & 静默丢弃，不 panic \\
  \texttt{SparseSystem::add\_diag}  & $i \ge \text{dim}$（越界）     & 静默丢弃，不 panic \\
  \texttt{conjugate\_gradient}      & $n = 0$（空系统）              & 返回 \texttt{Some(Vec::new())} \\
  \texttt{conjugate\_gradient}      & $\|b\| < 10^{-30}$             & 返回零解 \texttt{vec![0.0; n]} \\
  \texttt{conjugate\_gradient}      & $A$ 与 $b$ 维度不匹配           & \texttt{assert\_eq!} 触发 panic（编程错误） \\
  \texttt{conjugate\_gradient}      & \texttt{max\_iter} 未收敛       & 返回 \texttt{None} \\
  \texttt{conjugate\_gradient}      & $A$ 半正定（非 SPD）            & 调用方需先 \texttt{regularize\_diagonal} 正则化 \\
  \bottomrule
\end{tabular}
\end{center}

\section{使用示例}

\begin{lstlisting}[language=Rust]
use halfedge::linalg::{SparseSystem, conjugate_gradient, regularize_diagonal};

// 1. 构建稀疏对称系统（拉普拉斯风格）
let mut sys = SparseSystem::new(3);
sys.add(0, 0,  2.0);   // 对角
sys.add(1, 1,  2.0);
sys.add(2, 2,  2.0);
sys.add(0, 1, -1.0);   // 非对角，自动写入 (1,0)
sys.add(1, 2, -1.0);   // 非对角，自动写入 (2,1)
let mut a = sys.finish();

// 2. 半正定系统先正则化（加 lambda * I 使其 SPD）
regularize_diagonal(&mut a, 1e-4);

// 3. 共轭梯度求解 Ax = b
let b = vec![1.0, 0.0, -1.0];
match conjugate_gradient(&a, &b, 1000, 1e-10) {
    Some(x) => println!("解 x = {:?}", x),
    None    => panic!("CG 未收敛，请增大 max_iter 或检查正定条件"),
}
\end{lstlisting}

\section{测试覆盖}

\begin{center}
\renewcommand{\arraystretch}{1.18}
\begin{tabular}{@{}ll@{}}
  \toprule
  \textbf{测试名} & \textbf{验证点} \\
  \midrule
  \texttt{test\_cg\_identity}            & $A = I$ 时 \texttt{CG} 一步收敛，$x = b$ \\
  \texttt{test\_cg\_small\_spd}          & $3\times3$ SPD 系统，验证 $Ax \approx b$ 至 $10^{-10}$ \\
  \texttt{test\_regularize\_diagonal}    & 对角元正确累加 $\lambda$，非对角元不变 \\
  \texttt{test\_sparse\_system\_builder} & \texttt{add} 对称写入、\texttt{add\_diag} 对角写入正确 \\
  \bottomrule
\end{tabular}
\end{center}

\texttt{src/linalg.rs} 共 \textbf{4 个}单元测试，全库共 \textbf{371 个}。

\end{document}
